\begin{thebibliography}{99}
\bibitem{ref1} Ziv Bar-Yossef, T. S. Jayram, Ravi Kumar, and D. Sivakumar. An Information Statistics Approach to Data Stream and Communication Complexity. In Proceedings of the 43th Annual IEEE Symposium on Foundations of Computer Science, FOCS ’02, pages 209–218, 2002. doi:10.1109/SFCS.2002.1181944.
\bibitem{ref2} Boaz Barak, Mark Braverman, Xi Chen, and Anup Rao. How to Compress Interactive Com- munication. SIAM Journal on Computing, 42(3):1327–1363, 2013. doi:10.1137/100811969.
\bibitem{ref3} Mohammad Bavarian, Thomas Vidick, and Henry Yuen. Anchoring Games for Parallel Repetition, 2015. arXiv:1509.07466.
\bibitem{ref4} Mohammad Bavarian, Thomas Vidick, and Henry Yuen. Hardness Amplification for Entangled Games via Anchoring. In Proceedings of the 49th Annual ACM SIGACT Symposium on Theory of Computing, STOC ’17, page 303–316, 2017. doi:10.1145/3055399.3055433.
\bibitem{ref5} Avraham Ben-Aroya, Oded Regev, and Ronald de Wolf. A Hypercontractive Inequality for Matrix-Valued Functions with Applications to Quantum Computing and LDCs. In Proceedings of the 49th Annual IEEE Symposium on Foundations of Computer Science, FOCS ’08, pages 477–486, 2008. doi:10.1109/FOCS.2008.45.
\bibitem{ref6} Mario Berta, Matthias Christandl, and Renato Renner. The Quantum Reverse Shannon Theorem Based on One-Shot Information Theory. Communications in Mathematical Physics, 306(3):579–615, 2011. doi:10.1007/s00220-011-1309-7.
\bibitem{ref7} Mark Braverman. Interactive information complexity. SIAM Journal on Computing, 44(6):1698– 1739, 2015. doi:10.1137/130938517.
\bibitem{ref8} Mark Braverman, Ankit Garg, Young Kun Ko, Jieming Mao, and Dave Touchette. Near- Optimal Bounds on Bounded-Round Quantum Communication Complexity of Disjointness. In 2015 IEEE 56th Annual Symposium on Foundations of Computer Science, pages 773–791, 2015. doi:10.1109/FOCS.2015.53.
\bibitem{ref9} Mark Braverman and Gillat Kol. Interactive Compression to External Information. In Proceedings of the 50th Annual ACM SIGACT Symposium on Theory of Computing, STOC ’18, page 964–977, 2018. doi:10.1145/3188745.3188956.
\bibitem{ref10} Mark Braverman and Anup Rao. Information Equals Amortized Communication. IEEE Trans- actions on Information Theory, 60(10):6058–6069, 2014. doi:10.1109/TIT.2014.2347282.
\bibitem{ref11} Mark Braverman, Anup Rao, Omri Weinstein, and Amir Yehudayoff. Direct Product via Round-Preserving Compression. In Automata, Languages, and Programming, volume 7965 of Lecture Notes in Computer Science, pages 232–243. Springer Berlin Heidelberg, 2013. doi:10.1007/978-3-642-39206-1\_20.
\bibitem{ref12} Mark Braverman, Anup Rao, Omri Weinstein, and Amir Yehudayoff. Direct Products in Communication Complexity. In Proceedings of the 54th Annual IEEE Symposium on Foundations of Computer Science, FOCS ’13, pages 746–755, 2013. doi:10.1109/FOCS.2013. 85.
\bibitem{ref13} Mark Braverman and Omri Weinstein. An Interactive Information Odometer and Applications. In Proceedings of the Forty-Seventh Annual ACM Symposium on Theory of Computing, STOC ’15, page 341–350, 2015. doi:10.1145/2746539.2746548.
\bibitem{ref14} Amit Chakrabarti, Yaoyun Shi, Anthony Wirth, and Andrew Yao. Informational Complexity and the Direct Sum Problem for Simultaneous Message Complexity. In Proceedings of the 42nd Annual IEEE Symposium on Foundations of Computer Science, FOCS ’01, pages 270–278, 2001. doi:10.1109/SFCS.2001.959901.
\bibitem{ref15} Richard Cleve, William Slofstra, Falk Unger, and Sarvagya Upadhyay. Perfect Parallel Repetition Theorem for Quantum XOR Proof Systems. Computational Complexity, 17(2):282– 299, 2008. doi:10.1007/s00037-008-0250-4.
\bibitem{ref16} Irit Dinur. The PCP Theorem by Gap Amplification. J. ACM, 54(3):12–es, 2007. doi: 10.1145/1236457.1236459.
\bibitem{ref17} Irit Dinur, David Steurer, and Thomas Vidick. A Parallel Repetition Theorem for En- tangled Projection Games. Computational Complexity, 24(2):201–254, 2015. doi:10.1007/ s00037-015-0098-3.
\bibitem{ref18} Prahladh Harsha, Rahul Jain, David McAllester, and Jaikumar Radhakrishnan. The Commu- nication Complexity of Correlation. IEEE Transactions on Information Theory, 56(1):438–449, 2010.
\bibitem{ref19} Thomas Holenstein. Parallel Repetition: Simplifications and the No-Signaling Case. In Proceedings of the Thirty-Ninth Annual ACM Symposium on Theory of Computing, STOC ’07, page 411–419, 2007. doi:10.1145/1250790.1250852.
\bibitem{ref20} Rahul Jain. New Strong Direct Product Results in Communication Complexity. Journal of the ACM, 62(3), 2015. doi:10.1145/2699432.
\bibitem{ref21} Rahul Jain and Hartmut Klauck. New Results in the Simultaneous Message Passing Model via Information Theoretic Techniques. In Proceedings of the 24th Annual IEEE Conference on Computational Complexity, CCC ’09, pages 369–378, 2009. doi:10.1109/CCC.2009.28.
\bibitem{ref22} Rahul Jain, Hartmut Klauck, and Ashwin Nayak. Direct Product Theorems for Classical Communication Complexity via Subdistribution Bounds: Extended Abstract. In Proceedings of the 40th Annual ACM Symposium on Theory of Computing, STOC ’08, pages 599–608, 2008. doi:10.1145/1374376.1374462.
\bibitem{ref23} Rahul Jain and Ashwin Nayak. Short Proofs of the Quantum Substate Theorem. IEEE Transactions on Information Theory, 58(6):3664–3669, 2012.
\bibitem{ref24} Rahul Jain, Attila Pereszlényi, and Penghui Yao. A Parallel Repetition Theorem for Entangled Two-Player One-Round Games under Product Distributions. In 2014 IEEE 29th Conference on Computational Complexity (CCC ’14), pages 209–216, 2014.
\bibitem{ref25} Rahul Jain, Attila Pereszlényi, and Penghui Yao. A Direct Product Theorem for Two-Party Bounded-Round Public-Coin Communication Complexity. Algorithmica, 76(3):720–748, 2016. doi:10.1007/s00453-015-0100-0.
\bibitem{ref26} Rahul Jain, Jaikumar Radhakrishnan, and Pranab Sen. The Quantum Communication Complexity of the Pointer Chasing Problem: The Bit Version. In FSTTCS 2002: Foundations of Software Technology and Theoretical Computer Science, volume 2556 of Lecture Notes in Computer Science, pages 218–229, 2002. doi:10.1007/3-540-36206-1\_20.
\bibitem{ref27} Rahul Jain, Jaikumar Radhakrishnan, and Pranab Sen. A Direct Sum Theorem in Com- munication Complexity via Message Compression. In Automata, Languages and Program- ming, volume 2719 of Lecture Notes in Computer Science, pages 300–315. Springer, 2003. doi:10.1007/3-540-45061-0\_26.
\bibitem{ref28} Rahul Jain, Jaikumar Radhakrishnan, and Pranab Sen. A Lower Bound for the Bounded Round Quantum Communication Complexity of Set Disjointness. In Proceedings of the 44th Annual IEEE Symposium on Foundations of Computer Science, FOCS ’03, pages 220–229. IEEE Computer Society, 2003. doi:10.1109/SFCS.2003.1238196.
\bibitem{ref29} Rahul Jain, Jaikumar Radhakrishnan, and Pranab Sen. Prior Entanglement, Message Com- pression and Privacy in Quantum Communication. In 20th Annual IEEE Conference on Computational Complexity (CCC ’05), pages 285–296, 2005.
\bibitem{ref30} Rahul Jain, Jaikumar Radhakrishnan, and Pranab Sen. Optimal Direct Sum and Privacy Trade- off Results for Quantum and Classical Communication Complexity, 2008. arXiv:0807.1267.
\bibitem{ref31} Rahul Jain, Jaikumar Radhakrishnan, and Pranab Sen. A Property of Quantum Relative Entropy with an Application to Privacy in Quantum Communication. Journal of the ACM, 56(6), 2009. doi:10.1145/1568318.1568323.
\bibitem{ref32} Rahul Jain and Penghui Yao. A Strong Direct Product Theorem in Terms of the Smooth Rectangle Bound, 2012. arXiv:1209.0263.
\bibitem{ref33} Zhengfeng Ji, Anand Natarajan, Thomas Vidick, John Wright, and Henry Yuen. MIP*=RE, 2020. arXiv:2001.04383.
\bibitem{ref34} Julia Kempe, Oded Regev, and Ben Toner. Unique Games with Entangled Provers are Easy. SIAM Journal on Computing, 39(7):3207–3229, 2010. doi:10.1137/090772885.
\bibitem{ref35} Hartmut Klauck. A Strong Direct Product Theorem for Disjointness. In Proceedings of the 42nd ACM Symposium on Theory of Computing, STOC ’10, pages 77–86, 2010. doi: 10.1145/1806689.1806702.
\bibitem{ref36} Hartmut Klauck, Robert Špalek, and Ronald de Wolf. Quantum and Classical Strong Direct Product Theorems and Optimal Time-Space Tradeoffs. SIAM Journal on Computing, 36(5):1472–1493, 2007. doi:10.1137/05063235X.
\bibitem{ref37} Gillat Kol. Interactive Compression for Product Distributions. In Proceedings of the Forty- Eighth Annual ACM Symposium on Theory of Computing, STOC ’16, page 987–998, 2016. doi:10.1145/2897518.2897537.
\bibitem{ref38} Troy Lee, Adi Shraibman, and Robert Špalek. A Direct Product Theorem for Discrepancy. In Proceedings of the 23rd Annual IEEE Conference on Computational Complexity, CCC ’08, pages 71–80, 2008. doi:10.1109/CCC.2008.25.
\bibitem{ref39} Ran Raz. A Parallel Repetition Theorem. In Proceedings of the Twenty-Seventh Annual ACM Symposium on Theory of Computing, page 447–456, 1995. doi:10.1145/225058.225181.
\bibitem{ref40} Alexander A. Razborov. On the Distributional Complexity of Disjointness. Theoretical Computer Science, 106(2):385–390, 1992. doi:10.1016/0304-3975(92)90260-M.
\bibitem{ref41} Ronen Shaltiel. Towards Proving Strong direct Product Theorems. Computational Complexity, 12(1-2):1–22, 2003. doi:10.1007/s00037-003-0175-x.
\bibitem{ref42} Alexander A. Sherstov. Strong Direct Product Theorems for Quantum Communication and Query Complexity. SIAM Journal on Computing, 41(5):1122–1165, 2012. doi:10.1137/ 110842661.
\bibitem{ref43} Alexander A. Sherstov. Compressing Interactive Communication Under Product Distributions. SIAM Journal on Computing, 47(2):367–419, 2018. doi:10.1137/16M109380X.
\bibitem{ref44} Emanuele Viola and Avi Wigderson. Norms, XOR Lemmas, and Lower Bounds for Polynomials and Protocols. Theory of Computing, 4(7):137–168, 2008. doi:10.4086/toc.2008.v004a007.
\bibitem{ref45} Andrew Chi-Chin Yao. Probabilistic computations: Toward a unified measure of complexity. In 18th Annual Symposium on Foundations of Computer Science (SFCS 1977), pages 222–227, 1977. doi:10.1109/SFCS.1977.24.
\bibitem{ref46} Henry Yuen. A Parallel Repetition Theorem for All Entangled Games. In 43rd International Colloquium on Automata, Languages, and Programming (ICALP ’16), volume 55 of Leibniz International Proceedings in Informatics (LIPIcs), pages 77:1–77:13, 2016. doi:10.4230/ LIPIcs.ICALP.2016.77.
\end{thebibliography}
